\section{Proofs for the feasible-domain extensions}\label{app:con}

Throughout, $N_{\mathfrak B}$ denotes the number of bags of the tree
decomposition used by the algorithm, and $C_0$ an absolute constant that may
change between occurrences.

\subsection{Global graphs}\label{app:con:graph}

\begin{proof}[Proof of \cref{lem:con:expand}]
Substitute only explicit rational constants and implicit coordinates with
rational singleton brackets. A constant
implicit root need not be rational: the root of $y^2=2$ in $[1,2]$ has no
parents. Such coordinates remain in the chart and in the rational
original-domain formula used by the fallback. Their ancestor sets are empty
and do not affect scope containment or running intersection. In case (E),
induction along the dependency order shows that $\psi_j$ depends only on
$\Lambda(y_j)$ and that $\abs{\Lambda(y_j)}\le k^\delta$ for a coordinate at
dependency depth $\delta\le D_{\rm g}$: $\Lambda(y_j)$ is the union of the sets
$\Lambda(\pi)$ over its at most $k$ parents $\pi$. In case (I), $\psi_j$
depends on $x_{S_j}$ and $\abs{S_j}\le k$. Hence
$\abs{\beta'}\le\sum_{v\in\beta}\abs{\Lambda(v)}\le p\max\{1,k^{D_{\rm g}}\}$,
respectively $p\max\{1,k\}$.

Every retained coordinate lies in a bag of the original decomposition and
hence in the corresponding expanded bag. A reduced factor $f_a(x,\psi(x))$
depends only on $\bigcup_{v\in S_a}\Lambda(v)$, which lies in $\beta'$ for a
bag $\beta\supseteq S_a$. The term $\eta_j\psi_j$ depends on $\Lambda(y_j)$,
which lies in $\beta'$ for every bag $\beta$ containing $y_j$.

For running intersection, fix a retained coordinate $x_i$ and let $\mathcal
D_i$ be the set of dependent coordinates $y$ with $x_i\in\Lambda(y)$. The
expanded bags containing $x_i$ are those of the original bags that contain
$x_i$ or a member of $\mathcal D_i$, that is, the union of the subtrees
$\mathcal T_v$ of bags containing $v$, over $v\in\{x_i\}\cup\mathcal D_i$.
Each $\mathcal T_v$ is connected. In case (I), $y_j\in\mathcal D_i$ means
$i\in S_j$, and the bag containing $S_j\cup\{y_j\}$ lies in $\mathcal T_{x_i}$
and in $\mathcal T_{y_j}$. In case (E), a member of $\mathcal D_i$ is the end
of a dependency path starting at $x_i$, consecutive elements of which are a
parent and its child; the bag containing the child's constraint scope lies in
both of their subtrees. In both cases every $\mathcal T_y$, $y\in\mathcal D_i$,
is connected to $\mathcal T_{x_i}$ within the union, which is therefore
connected.
\end{proof}

\begin{proof}[Proof of \cref{thm:con:graph}, case (E)]
\emph{Reduction.} Substitution in the dependency order computes every
$\psi_j$ as an explicit polynomial of degree at most $e^{D_{\rm g}}$ in at
most $\min\{n,k^{D_{\rm g}}\}$ retained variables, so it has
$\poly(I)$ monomials, and similarly for the reduced factors, whose degree is
at most $d'=de^{D_{\rm g}}$; the base part $F_0(x,\psi(x))$ is computed once,
and the sampled part $\eta^{\mathsf T}\psi$ adds coefficients of length
$\poly(I+b)$. The production cost is also $\poly(I+b)$, with exponents
depending only on $d$, $e$ and $D_{\rm g}$, as required by
\cref{lem:con:uniform}. The map $(x,y)\mapsto x$ is a bijection from $\mathfrak X$ onto
$X$ with inverse $x\mapsto(x,\psi(x))$, and
$F_{\gamma,\eta}(x,\psi(x))=G_\eta(x)+\gamma^{\mathsf T}x$. Hence minimizers
correspond, and the sampled problem is that of \cref{lem:con:uniform} for the
family $\{G_\eta\}$ with the expanded decomposition of \cref{lem:con:expand}.

\emph{Hypotheses of \cref{lem:con:uniform}.} The bound $L$ of
\eqref{eq:con:curv} is valid for every $\eta$ in the noise box. Monomial
bounds for the derivatives of $F_0(x,\psi(x))$ and of the $\psi_j$ on
$\widehat X$, the latter multiplied by $\sigma\ge\abs{\eta_j}$, give uniform
$\kappa_2$ and $\kappa_3$ of polynomial encoding length. For~(b), apply
\cref{thm:count:fallback} to the base instance whose domain formula consists
of the bounds of $X$, the integer constraints, and the equations
$y_j-P_j(x,y_1,\dots,y_{j-1})=0$, with objective $F_0$ tilted by
$(\gamma,\eta)$. It returns the lexicographically least minimizer
$(x^\circ,y^\circ)$, and $x^\circ$ minimizes $G_\eta+\gamma^{\mathsf T}x$ on
$X$. Its factor $B$ depends only on the base instance. \Cref{lem:con:uniform}
gives (i) and (ii); the charted patch output consists of the patch of the
reduced problem and the polynomials $P_j$.

\emph{Evaluation.} If the retained patch is a point, substitute it into the
explicit chart and evaluate the sampled polynomial exactly; no convex
optimization is needed. For a positive-dimensional patch, let
$K_\psi=1+\sum_jG_{\psi_j}$, where $G_{\psi_j}$ is a
rational monomial bound for $\sup_{\widehat X}\norm{\nabla\psi_j}_1$. For
$x,x'\in\widehat X$,
$\norm{(x,\psi(x))-(x',\psi(x'))}_2\le\norm{x-x'}_2+\sum_j\abs{\psi_j(x)-\psi_j(x')}
\le K_\psi\norm{x-x'}_2$. On ordinary draws, \cref{prop:sp:eval}(a) applied
to the reduced polynomial gives a rational $\tilde x$ in the patch with
$\norm{\tilde x-\hat x}_2\le2^{-q}/K_\psi$ and a value interval of width at
most $2^{-q}$; the lift $(\tilde x,\psi(\tilde x))$ is rational, satisfies the
equations, and lies within $2^{-q}$ of the optimizer, and its value
$G_\eta(\tilde x)+\gamma^{\mathsf T}\tilde x$ is the upper endpoint. On
fallback draws, \cref{prop:sp:eval}(b) applied to the retained coordinates of
$(x^\circ,y^\circ)$, with the chart map, gives the same. The extra precision
$\log K_\psi$ is polynomial in $I$.
\end{proof}

\begin{proof}[Proof of \cref{lem:con:chart}]
\emph{Roots.} For $x\in\widehat X_{S_j}$, the function $y\mapsto q_j(x,y)$ is
strictly increasing on $V_j$ and has a sign change there by
\eqref{eq:con:bracket}, so it has exactly one root $\psi_j(x)$ in $V_j$. If
$\ell'_j=u'_j$, then $q_j(x,\ell'_j)=0$ for all $x$, and the coordinate is
substituted.

\emph{Smoothness.} Fix $x^0\in\widehat X_{S_j}$ and $y^0=\psi_j(x^0)$. Since
$\partial_yq_j(x^0,y^0)\ge\mu_j>0$, the implicit function theorem gives an
open neighborhood $\mathcal N$ of $x^0$, a number $\epsilon>0$, and an
infinitely differentiable $\varphi$ on $\mathcal N$ such that $\varphi(x)$ is
the only root of $q_j(x,\cdot)$ in $(y^0-\epsilon,y^0+\epsilon)$ for
$x\in\mathcal N$. For $y\in V_j$ with $\abs{y-y^0}\ge\epsilon$, monotonicity
gives $\abs{q_j(x^0,y)}\ge\mu_j\epsilon$, so by uniform continuity
$q_j(x,y)\ne0$ for all such $y$ and all $x\in\widehat X_{S_j}$ close to $x^0$.
For these $x$, $\psi_j(x)$ lies in $(y^0-\epsilon,y^0+\epsilon)$ and therefore
equals $\varphi(x)$. Thus $\psi_j$ agrees near every point of
$\widehat X_{S_j}$ with an infinitely differentiable function, and its
derivatives are obtained by differentiating $q_j(x,\psi_j(x))=0$.

\emph{Derivative bounds.} Write $q$ for $q_j$, subscripts for partial
derivatives, and evaluate at $(x,\psi(x))$, where $q_y\ge\mu_j$ and all
partial derivatives of orders one to three are bounded by $A_j$. Then
$\psi_i=-q_i/q_y$, so $\abs{\psi_i}\le\Pi_{1,j}$, and
\[
 \psi_{ik}=-\frac{N_{ik}}{q_y},\qquad
 N_{ik}=q_{ik}+q_{iy}\psi_k+q_{ky}\psi_i+q_{yy}\psi_i\psi_k ,
\]
so $\abs{\psi_{ik}}\le A_j(1+\Pi_{1,j})^2/\mu_j=\Pi_{2,j}$. Differentiating
$q_y\psi_{ik}=-N_{ik}$ with respect to $x_l$ gives
$q_y\psi_{ikl}=-\frac{d}{dx_l}N_{ik}-\psi_{ik}\frac d{dx_l}q_y$. Here
$\abs{\frac d{dx_l}q_y}=\abs{q_{yl}+q_{yy}\psi_l}\le A_j(1+\Pi_{1,j})$, and the
total derivative of $N_{ik}$ is a sum of the terms $q_{ikl}+q_{iky}\psi_l$,
$(q_{iyl}+q_{iyy}\psi_l)\psi_k+q_{iy}\psi_{kl}$, the analogous term with $i$
and $k$ exchanged, and
$(q_{yyl}+q_{yyy}\psi_l)\psi_i\psi_k+q_{yy}(\psi_{il}\psi_k+\psi_i\psi_{kl})$.
Bounding each factor gives
$A_j[(1+\Pi_{1,j})^3+2(1+\Pi_{1,j})\Pi_{2,j}]$, and adding
$\Pi_{2,j}A_j(1+\Pi_{1,j})$ and dividing by $\mu_j$ gives $\Pi_{3,j}$.

\emph{Oracle.} If there are no dependent coordinates, evaluate the rational
polynomial and its derivatives directly. Otherwise, at a rational $x$,
bisection on $V_j$ with exact rational
signs of $q_j(x,\cdot)$ keeps a bracket $[l_j,u_j]\subseteq V_j$ with
$q_j(x,l_j)\le0\le q_j(x,u_j)$; monotonicity shows that it contains
$\psi_j(x)$, and a zero sign identifies the root exactly. Width $\epsilon$
takes $O(\log((u'_j-\ell'_j)/\epsilon))$ steps, each a fixed-degree
evaluation at a rational point whose length grows by one bit per step. By
the chain rule and the formulas above, $G_\eta(x)$ and every entry of
$\nabla G_\eta(x)$ and $\nabla^2G_\eta(x)$ equal $\Phi(x,\psi(x))$ for an
explicit expression $\Phi(x,y)$: a polynomial in $\eta$ and in partial
derivatives of $F_0$ and of the $q_j$ at $(x,y)$, divided by powers of the
$\partial_yq_j(x,y_j)$. On $\widehat X\times\prod_jV_j$ these denominators
are at least $\mu_j$, so $\abs{\partial_{y_j}\Phi}$ is bounded by a rational
$\Lambda_\Phi$ of polynomial encoding length, obtained from monomial bounds.
With brackets of widths $\epsilon_j$ and their rational midpoints $\tilde y$,
$\abs{\Phi(x,\psi(x))-\Phi(x,\tilde y)}\le\Lambda_\Phi\sum_j\epsilon_j/2$,
because the segment from $\psi(x)$ to $\tilde y$ stays in $\prod_jV_j$.
Choosing $\epsilon_j\le\epsilon/(m\Lambda_\Phi)$ gives the enclosures, and
$\Phi(x,\tilde y)$ is an exact rational evaluation. The number of
expressions is polynomial, and so is the cost. The same chain-rule
expressions and the bounds $\Pi_{\nu,j}$ give $\kappa_2$, $\kappa_3$ and
$K_\psi$, uniformly for $\abs{\eta_j}\le\sigma$.
\end{proof}

\begin{proof}[Proof of \cref{lem:con:approx}]
Write $G=G_\eta+\gamma^{\mathsf T}x$ and, for an allowed grid point $y$,
$\underline H(y)=\sum_\beta\underline\phi_\beta(y_\beta)$. Each reduced term
is assigned to exactly one bag, so
$G(y)-D_j\le\underline H(y)\le G(y)$ with $D_j=N_{\mathfrak B}\delta_j=E_j$.
\Cref{lem:sp:round} applies to $G$, whose coordinate curvature is at most $L$
on $\widehat X$ by \eqref{eq:con:curv}. For $x\in\Omega_j$ in a candidate cell
$C$ of $\beta$, the rounded point $Y$ satisfies
$G(Y)\ge\underline H(Y)\ge\underline q_C$, so $\underline q_C-E_j\le G(x)$.
The dynamic-programming minimizer $y^{(j)}$ has
$G(y^{(j)})\le\underline m_j+D_j\le\underline m_j+E_j$, so $U_j$ is an upper
bound for the value of a feasible point and $U_j\ge f^*$; rounding a global
minimizer gives $\underline m_j\le f^*+E_j$ and $U_j\le f^*+2E_j$. As in
\cref{prop:sp:prune}, every minimizer is retained and every removed point has
value above $U_j$. A retained cell has $\underline q_C\le U_j+E_j$, attained at
an allowed $y$ with $y_\beta$ a corner, and
$G(y)\le\underline H(y)+D_j\le f^*+4E_j$. Initialize $U_{-1}=+\infty$, and
store the point associated with the smallest certified upper value:
replace the stored point by $y^{(j)}$ when
$\underline m_j+E_j<U_{j-1}$, and otherwise keep the earlier point.
Its implicit lift is feasible and has value at most $U_j$; its exact
value is never compared. Update this pair before pruning.
\end{proof}

The certified rows at level $j$ need precision
$\delta_j=nLh_j^2/(8N_{\mathfrak B})$, of encoding length $\poly(I+j)$, so the
cost per row is polynomial by \cref{lem:con:chart}. Rows reused from an
earlier level are recomputed at the new precision. With the witness tolerance
$4E_j$, the proof of \cref{prop:sp:count} applies to the reduced objective
conditionally on $\eta$, with interval lengths
$La_i+8E_j/a_i\le(1+n)La_i$; the expected number of retained cells of a bag
is at most $2^{\abs\beta}\prod_{i\in\beta}[4+(1+n)Lw_i/(2\sigma)]$.

\begin{lemma}[Certified closure]\label{lem:con:approx-closure}
In case (I), the closure test with certified derivatives described in
\cref{sec:con:graph} is sound on every draw. If $g_*\ge g_0$, the unique
minimizer $a$ has $\abs{\partial_iG(a)}>\tau$ at every active continuous
coordinate, $\epsilon_g\le\tau/8$, $\epsilon_H\le g_0/8$, $A=2+nL/g_0$ and
$h_j\le\min\{1/(4A),\tau/(4\kappa_2A),g_0/(4\kappa_3A)\}$, then the test of
level $j$ succeeds.
\end{lemma}

\begin{proof}
For $x$ in the hull box, $\abs{\partial_iG(x)-\partial_iG(c)}\le\kappa_2r$ and
$\abs{\hat g_i-\partial_iG(c)}\le\epsilon_g$, so the enclosure
$\hat g_i\pm(\kappa_2r+\epsilon_g)$ contains $\partial_iG(x)$, and the
argument of \cref{prop:sp:closure-sound} applies. On the patch,
$\nabla^2_{\bar S\bar S}G\succeq\nabla^2_{\bar S\bar S}G(\bar x_S,c^P)-\kappa_3r^PI
\succeq\hat H-(\epsilon_H+\kappa_3r^P)I\succ g_0I$. For success, the witnesses
satisfy $g_0\norm{y-a}^2\le4E_j$, so $\norm{y-a}_2\le h_j\sqrt{nL/(2g_0)}$ and
every hull lies within $h_j(1+\sqrt{nL/(2g_0)})\le Ah_j$ of $a$. As in
\cref{prop:sp:closure-stop}, integer hulls are the correct singletons, and an
active coordinate with $\partial_iG(a)>\tau$ has enclosure lower end at least
$\partial_iG(a)-2\kappa_2r-2\epsilon_g\ge\tau-\tau/2-\tau/4>0$. The free
Hessian at $a$ is at least $2g_0I$, so
$\hat H-(\kappa_3r^P+\epsilon_H+g_0)I\succeq
(g_0-2\kappa_3r^P-2\epsilon_H)I\succeq\frac{g_0}4I$.
\end{proof}

\begin{lemma}[Active retained gradients on implicit graphs]\label{lem:con:kkt}
In case (I), let $D_{\rm I}=\max\{2,d,e\}$ and
$K_{\rm I}=\max\{1,n_c3^{n_c}N_ZD_{\rm I}^{n_c+2m}\}$. Conditionally on every
$\eta$ and for every $\tau>0$, the probability that $g_*>0$ and the unique
minimizer $a$ of $G=G_\eta+\gamma^{\mathsf T}x$ has a continuous coordinate
at a bound with $\abs{\partial_iG(a)}\le\tau$ is at most
$K_{\rm I}(\tau/\sigma+1/M)$.
\end{lemma}

\begin{proof}
Fix an integer assignment $z$, a face of the continuous retained box given by
a partition $I_C=\mathsf L\cup\mathsf U\cup\mathsf F$, and an index
$i\in\mathsf L\cup\mathsf U$; condition on $\eta$ and on $\gamma_l$, $l\ne i$.
Let $\tilde x(u)$ be the retained point with integer part $z$, bounds on
$\mathsf L\cup\mathsf U$, and $u\in\R^{\mathsf F}$ on the free set; put
$k=\abs{\mathsf F}$ and $\mathcal L=F_0(\tilde x(u),y)+\gamma^{\mathsf
T}\tilde x(u)+\eta^{\mathsf T}y+\sum_j\lambda_jq_j(\tilde x(u)_{S_j},y_j)$.
The system
\[
 q_j(\tilde x(u)_{S_j},y_j)=0,\qquad \partial_{y_j}\mathcal L=0,\qquad
 \partial_{u_l}\mathcal L=0\qquad(j\le m,\ l\in\mathsf F)
\]
has $k+2m$ equations in the unknowns $(u,y,\lambda)$, each of degree at most
$D_{\rm I}$, and does not involve $\gamma_i$. If $k+2m=0$, its unique
empty tuple is counted with weight one. Otherwise by \cref{lem:sp:bezout} it has at
most $D_{\rm I}^{k+2m}$ nonsingular zeros. For a zero $\varrho=(u,y,\lambda)$ put
$b_\varrho=\partial_{x_i}F_0(\tilde x(u),y)+\sum_j\lambda_j\partial_{x_i}q_j(\tilde
x(u)_{S_j},y_j)$, a complex number fixed by the conditioning. The set of real
$\gamma_i$ with $\abs{\gamma_i+b_\varrho}\le\tau$ lies in an interval of length
$2\tau$, so it has conditional probability at most $\tau/\sigma+1/M$. At a zero
that comes from a graph point, $\partial_yq_j\ge\mu_j>0$, the multipliers are
$\lambda_j=-(\partial_{y_j}F_0+\eta_j)/\partial_yq_j$, and the chain rule with
$\partial_i\psi_j=-\partial_{x_i}q_j/\partial_yq_j$ shows that the reduced
derivative there is $\partial_iG=\gamma_i+b_\varrho$.

Suppose $g_*>0$ and let $a$ be the minimizer, with its integer part, bounds
and free set as above. Then $(u,y,\lambda)=(a_{\mathsf F},\psi(a),\lambda(a))$
is a zero, because the reduced gradient on $\mathsf F$ vanishes. Write
$C=[q_u\ q_y]$ for the constraint Jacobian, with $q_y$ diagonal and
invertible, $Z=[I;-q_y^{-1}q_u]$, and $\mathsf H=\nabla^2_{(u,y)}\mathcal L$.
Since $\mathcal L(u,\psi(u),\lambda)$ equals the reduced objective for fixed
$\lambda$ and $\partial_y\mathcal L=0$ at the zero, the chain rule gives
$Z^{\mathsf T}\mathsf HZ=\nabla^2_{\mathsf F\mathsf F}G(a)\succeq2g_*I$, by
two-sided growth along the free directions. If the Jacobian of the system,
$\bigl[\begin{smallmatrix}\mathsf H&C^{\mathsf T}\\C&0\end{smallmatrix}\bigr]$,
annihilates $(v,\mu)$, then $Cv=0$ gives $v=Zw$, and multiplying the first
block row by $Z^{\mathsf T}$ gives $Z^{\mathsf T}\mathsf HZw=0$, so $w=0$, $v=0$,
and $q_y^{\mathsf T}\mu=0$ gives $\mu=0$. The zero is therefore nonsingular,
and the event is contained in the union of the events counted above, over at
most $N_Z3^{n_c}n_c$ choices.
\end{proof}

\begin{proof}[Proof of \cref{thm:con:graph}, case (I)]
\emph{Schedule.} Let $B$ be the factor of \cref{thm:count:fallback} for the
base instance with domain formula consisting of the bounds and integer
constraints of $X$, the bounds $y_j\in V_j$, and the equations $q_j=0$, and
objective $F_0$ tilted on all coordinates. Let $C_{\rm tail}$ be the constant
of \cref{lem:count:finite-tails} for this formula, with $n$ noisy variables
$x$, $m$ further variables $y$, and degree $\max\{d,e\}$; and $K_{\rm I}$ as in
\cref{lem:con:kkt}. Put $\rho=1/(4B)$, $g_0=\rho\sigma/(2W)$,
$\tau=\rho\sigma/(2K_{\rm I})$, $A=2+nL/g_0$, $\epsilon_g=\tau/8$,
$\epsilon_H=g_0/8$, let $J$ be the least level satisfying the condition of
\cref{lem:con:approx-closure}, and let $M$ be the least power of two with
$M\ge\max\{2,2^J,4nC_{\rm tail}/\rho,2K_{\rm I}/\rho\}$. All these have
polynomial encoding length, and $\kappa_2$, $\kappa_3$ are uniform in $\eta$
(\cref{lem:con:chart}).

\emph{Tails.} Conditionally on $\eta$, the set $\mathfrak X$ is compact and
described by the formula above, $F=F_0+\eta^{\mathsf T}y$ has rational
coefficients, and the reduced objective $\min_y\{F(x,y):(x,y)\in\mathfrak X\}$
is $G_\eta(x)$, because the fiber is the single point $\psi(x)$. Hence
\cref{lem:count:finite-tails}(a),(b) give
$\Prob\{g_*<g_0\mid\eta\}\le g_0W/\sigma+2nC_{\rm tail}/M\le\rho$, and
\cref{lem:con:kkt} gives probability at most $\rho$ for small active
gradients. By \cref{lem:con:approx-closure}, closure fails at level $J$ with
probability at most $2\rho=1/(2B)$, conditionally on every $\eta$.

\emph{Correctness and work.} Closure outputs are correct by
\cref{lem:con:approx,lem:con:approx-closure} and the strong convexity of the
reduced patch. Otherwise \cref{thm:count:fallback} is applied to
$F_{\gamma,\eta}$ on $\mathfrak X$. The work of the levels is bounded as in the
proof of \cref{thm:sp:main}, with the count above and the polynomial cost per
certified row; the fallback contributes $\poly(I)$ in expectation by
\cref{lem:count:rare-fallback}(a).

\emph{Evaluation.} First handle a point patch. Its retained coordinates are
rational and fixed, but its dependent coordinates can still be irrational.
Use \cref{lem:con:chart} to enclose the sampled objective to width $2^{-q}$
and to refine each dependent root to width at most
$2^{-q-1}/\max\{1,m\}$. Their exact roots give the feasible lift; their
rational midpoints give an ambient approximation within $2^{-q}/2$ of it.
This costs $\poly(I+q)$, with $q+O(\log(n+m+1))$ coordinate precision
and the oracle's polynomial encoded allowance for objective evaluation.
No positive-dimensional convex optimization is invoked.

For a positive-dimensional ordinary patch apply \cref{lem:sp:gls} to the reduced
objective on the patch box, with $Z=I$, coordinatewise clipping as repair,
$\nu=\norm\cdot_2$, chain-rule bounds $V$ and $G$, and the approximate oracle
of \cref{lem:con:chart}; its curvature hypothesis holds with $g=g_0$. With
$\eta_{\rm GLS}=\min\{2^{-q},\frac{g_0}2(2^{-q}/(2K_\psi))^2\}$ it returns a
rational $\tilde x$ with $\norm{\tilde x-\hat x}_2\le2^{-q}/(2K_\psi)$ and a
value interval of width at most $2^{-q}$. By \eqref{eq:con:Kpsi} the exactly
feasible lift $(\tilde x,\psi(\tilde x))$, given by $\tilde x$ and root
brackets, lies within $2^{-q}/2$ of the optimizer, and refining the brackets to
total width $2^{-q}/2$ gives a rational point within $2^{-q}$. On fallback
draws, let $G'\ge\max\{1,\sup_{\widehat X}\norm{\nabla G}_1\}$ be a rational
chain-rule bound, refine each retained continuous coordinate to error at
most $2^{-q}/(4n\max\{1,K_\psi,G'\})$, clip it to its original interval,
and identify the integer labels exactly as in \cref{prop:sp:eval}(b).
The retained distance is at most $2^{-q}/(4K_\psi)$ and the objective gap
at most $2^{-q}/4$. Attach the exact chart roots, refine their brackets to
total width $2^{-q}/2$, and refine the optimal value to width $2^{-q}/2$.
These give both the feasible lift with a certified objective gap and the
ambient rational approximation. All extra precision is $q+\poly(I)$ and
costs $B\poly(I+q)$ under \cref{thm:count:fallback}.
\end{proof}

\subsection{Affine-state recurrences}\label{app:con:actuator}

\begin{proof}[Proof of \cref{thm:con:actuator}]
\emph{Feasible set.} For every free point the recurrence defines one
trajectory, and by \eqref{eq:con:invariant} and induction on $k$,
$s_k\in I_k$ for all $k$, because $u_k\in\operatorname{hull}U_k$. Conversely
every feasible trajectory arises in this way. The identity
\eqref{eq:con:adjoint} follows by substituting $s_N=a_{N-1}s_{N-1}+\chi_{N-1}(u_{N-1})$
into $(c_N+\eta_N)s_N$, which turns the coefficient of $s_{N-1}$ into
$c_{N-1}+\eta_{N-1}+a_{N-1}\lambda_N=\lambda_{N-1}$, and repeating down to
$\lambda_1s_1=\lambda_1(a_0s_0+\chi_0(u_0))$. Hence, with constants from
substituted fixed coordinates carried separately, the sampled objective of a
trajectory equals $Q_\eta(z)+\xi^{\mathsf T}z$ at its free point.

Substitute fixed free coordinates before computing widths or a schedule.
If none remains, use $M=2$, draw the ambient coefficients once, and return
the unique rational trajectory and its exact sampled value. The remaining
argument concerns a nonempty free box with positive coordinate widths.

\emph{Uniform data.} By induction from $\abs{\lambda_N}\le C_s+\sigma$ and
$\abs{\lambda_k}\le C_s+\sigma+a\abs{\lambda_{k+1}}$, every
$\abs{\lambda_k}\le\Lambda$. The added terms are unary, so the coordinate
curvature, the Hessian row sums and the third-derivative row sums of
$Q_\eta$ exceed those of $q$ by at most $\Lambda G_2$, $\Lambda G_2$ and
$\Lambda G_3$, which gives \eqref{eq:con:Lact} and uniform $\kappa_2,\kappa_3$.
The adjoints are computed by $N$ affine steps, so their bit lengths are sums
of the lengths of the $a_k$, $c_k$ and $\eta_k$, polynomial in $I+b$.
Producing the coefficients of $Q_\eta$ also costs $\poly_d(I+b)$; neither
height nor production cost is asserted to be bounded by the literal $I+b$. The
fallback of \cref{lem:con:uniform}(b) is \cref{thm:count:fallback} applied to
the base instance with variables $(s,u)$, domain formula
$\bigwedge_k[s_{k+1}=a_ks_k+\chi_k(u_k)]$ together with the variable bounds and
integer constraints, and objective $F_0$ tilted on all coordinates; its
minimizers are the trajectories of the minimizers of $Q_\eta+\xi^{\mathsf
T}z$. \Cref{lem:con:uniform} with the free-variable decomposition gives (i)
and (ii).

\emph{Evaluation.} For a rational free point, the states are rational: each
step evaluates a fixed-degree polynomial at its own control and adds an affine
function of the previous state, so bit lengths add over the horizon instead
of being raised to powers, and the trajectory has length polynomial in $I$
and the length of the free point. It is exactly feasible, and integer labels
are exact. For two free points, $\abs{\Delta s_{k+1}}\le a\abs{\Delta
s_k}+G_1\abs{\Delta u_k}$, so the vector $(\abs{\Delta s_k})_{k\ge1}$ is
dominated by a lower-triangular Toeplitz matrix with entries $a^l$ applied to
$(\abs{\Delta s_0},G_1\abs{\Delta u_0},\dots)$. Its row and column sums are at
most $1/(1-a)$, so its spectral norm is at most $1/(1-a)$, and the trajectory
map has Lipschitz constant at most $K_{\rm traj}=1+(1+G_1)/(1-a)$. The free
evaluation of \cref{prop:sp:eval}, with $\log_2K_{\rm traj}$ extra bits,
gives (iii); objective values and gaps transfer exactly by
\eqref{eq:con:adjoint}.
\end{proof}

\subsection{Products of simplices}\label{app:con:simplex}

At levels with $h_j\ge1$ the cell of a block is its original simplex and
its corners and grid are the original simplex vertices: $0,e_i$ for an
inequality block and $e_i$ for an equality block. At levels with
$h=h_j=1/m<1$ the cells are the nonempty sets
$C=\Delta_b\cap\prod_{i\in b}[hk_i,h(k_i+1)]$, with
$k_i\in\{0,\dots,m-1\}$, or the same with $\Delta^=_b$; their corners are
their points in $(h\Z)^b$. These bounded indices exclude exterior cubes
that would otherwise add duplicate boundary cells. Integer coordinates use
the cells of \cref{sec:sp:cells}. A bag cell is a product of block cells
and integer cells.

\begin{lemma}[Feasible block rounding]\label{lem:con:simplex-round}
Every block cell is the convex hull of its corners, which are original
simplex vertices when $h_j\ge1$ and cube corners when $h_j<1$.
For $x\in\Omega_j$ and a cell $C_0\in\mathcal W_j(\beta_0)$ containing
$x_{\beta_0}$, there is a random $Y\in A_j$ with $Y_{\beta_0}$ a corner of
$C_0$, $Y_i=x_i$ whenever $x_i$ is a grid value, and
$\E F_\gamma(Y)\le F_\gamma(x)+nLh_j^2/8$.
\end{lemma}

\begin{proof}
For $h_j<1$, with $z_i=x_i/h-k_i$, a cell becomes
$\{z\in[0,1]^b:\sum_iz_i\le\rho\}$, or
the same with equality, where $\rho=m-\sum_ik_i$ is an integer. At a vertex,
$b$ linearly independent constraints are tight. If the sum constraint is not
among them, every coordinate is $0$ or $1$; if it is, at least $b-1$
coordinates are $0$ or $1$, and the remaining one equals $\rho$ minus an
integer and lies in $[0,1]$. So all vertices are cube corners, and the
bounded polytope is their convex hull. A whole simplex has the vertices $0$
and $e_i$, or $e_i$.

For each block $b$, choose a cell $C_b\ni x_b$ (the one of $C_0$ if $b\subseteq
\beta_0$), write $x_b$ as a convex combination of the corners of $C_b$, and let
$Y_b$ equal each corner with its coefficient, independently over blocks;
round integer coordinates as in \cref{lem:sp:round}. Then $\E Y_b=x_b$, and
each coordinate of $Y_b$ takes the two values $hk_i,h(k_i+1)$, so its variance
is at most $h^2/4$ (at most $1/4\le h^2/4$ while $h\ge1$). If $x_i=hk_i$, then
$Y_i\ge hk_i$ with mean $hk_i$ forces $Y_i=x_i$; the same holds at the upper
cube face. Hence $Y_b$ lies in every level-$j$ cell of $b$ that contains $x_b$:
coordinates of $x_b$ strictly inside a grid interval have the same interval in
every such cell, coordinates on a grid value stay fixed, and $Y_b$ is
feasible. As a grid point of that cell it is one of its corners. Every bag
cell containing $x_\beta$ therefore contains $Y_\beta$ as a corner, so
$Y\in A_j$. Rounding the blocks one after another, the segment from $x_b$ to
$Y_b$ lies in $C_b$, so \eqref{eq:con:blockcurv} and Taylor's theorem give
$\E[F_\gamma\mid\text{earlier blocks}]\le F_\gamma+\frac L2\E\norm{Y_b-x_b}^2$,
and summing over blocks and integer coordinates gives the allowance.
\end{proof}

With this lemma the proof of \cref{prop:sp:prune} applies verbatim.

\begin{lemma}[Count by faces]\label{lem:con:simplex-count}
Let $j\le J$ and $M\ge2^J$. For a bag $\beta$, the set $N_j(\beta)$ of feasible
bag grid tuples $v$ with $V_\beta(v)+\gamma_\beta^{\mathsf T}v\le f^*+2E_j$
satisfies
\[
 \E\abs{N_j(\beta)}\le4^{\abs\beta}\Bigl[4+\frac{(2+n/2)L\max\{1,w_{\max}\}}{2\sigma}\Bigr]^{\abs\beta},
\]
and $\abs{\mathcal W'_j(\beta)}\le2^{\abs\beta}\abs{N_j(\beta)}$ on every draw.
\end{lemma}

\begin{proof}
If $h_j\ge1$, each block has at most $\abs b+1\le2^{\abs b}$ grid vertices.
Use their trivial probability bound one and the integer-coordinate count
of \cref{prop:sp:count}; this is already bounded by the displayed product.
Thus the following face comparisons concern $h_j<1$ only.

Because bags contain whole blocks, the domain outside $\beta$ is a fixed
product, and $V_\beta$ is defined as in \eqref{eq:sp:V}. For a feasible bag
tuple $v$ and a direction $\mathsf d$ inside one block $b\subseteq\beta$ with
$v\pm h\mathsf d$ feasible, the minimizer argument of \cref{lem:sp:compare}
and \eqref{eq:con:blockcurv} give
$V_\beta(v+h\mathsf d)+V_\beta(v-h\mathsf d)-2V_\beta(v)\le Lh^2\norm{\mathsf d}^2$.

Fix $v$ and a block $b\subseteq\beta$, let $\mathsf S$ be the positive support
of $v_b$, and let $r_b$ be the dimension of the smallest face of the
simplex containing $v_b$. If the budget is slack, $r_b=\abs{\mathsf S}$, and
the directions $\mathsf d=e_i$, $i\in\mathsf S$, satisfy $v\pm h\mathsf d$
feasible, because positive coordinates and the budget slack are positive
multiples of $h$. If the budget is tight (always for equality blocks),
$r_b=\abs{\mathsf S}-1$; with the least index $\iota\in\mathsf S$ as anchor, the
directions $e_i-e_\iota$, $i\in\mathsf S\setminus\{\iota\}$, are feasible in
both signs. Condition on the noise outside $\beta$, on the anchor
coefficients of tight blocks, and on the coefficients of the zero coordinates.
If $v\in N_j(\beta)$, comparison with $v\pm h\mathsf d$ as in
\cref{lem:sp:compare} confines $\gamma^{\mathsf T}\mathsf d$, which is $\gamma_i$
or $\gamma_i$ minus a conditioned anchor, to an interval determined by the
conditioned values and $v$, of length at most
$Lh\norm{\mathsf d}^2+4E_j/h\le(2+n/2)Lh$. Each direction involves a different
remaining coefficient, so the conditional probability is at most
$[(2+n/2)Lh/(2\sigma)+1/M]^{r_b}$ for the block, times the integer factors of
\cref{prop:sp:count}.

A face of dimension $r$ has $\binom{1/h-1}{r}\le h^{-r}$ grid points in its
relative interior: positive multiples of $h$ in the slack case with total at
most $1-h$, and compositions of $1$ in the tight case. The simplex has at most
$2^{\abs b+1}$ faces, and $1/(Mh)=2^j/(Ms)\le1$. Hence the expected
contribution of the block is at most
$2^{\abs b+1}[1+(2+n/2)L/(2\sigma)]^{\abs b}\le4^{\abs b}[4+(2+n/2)L/(2\sigma)]^{\abs b}$;
while $h\ge1$ the block grid consists of at most $\abs b+1$ vertices. The
probabilities of different blocks and integer coordinates multiply after the
conditioning, and summing over tuples factorizes, which proves the bound.
Every retained cell has a corner in $N_j(\beta)$ by \cref{prop:sp:prune}(d),
and a grid tuple is a corner of at most $2^{\abs\beta}$ cells, since it lies in
at most $2^{\abs b}$ cubes per block.
\end{proof}

\begin{lemma}[Simplex closure]\label{lem:con:simplex-close}
Let $Q$ be the intersected hull box of level $j$, with midpoint $c$ and
largest half-width $r$, and let $\kappa_2,\kappa_3$ hold on the enclosing
coordinate box. The following recorded equalities hold at every global
minimizer:
\begin{enumerate}[label=(\alph*)]
\item for an inequality block, $x_i=0$ if $\partial_iF_\gamma(c)-\kappa_2r>0$;
\item for an inequality block, $\sum_{i\in b}x_i=1$ if
$\partial_lF_\gamma(c)+\kappa_2r<0$ for some $l\in b$;
\item for any block, $x_i=0$ if
$\partial_iF_\gamma(c)-\partial_lF_\gamma(c)-2\kappa_2r>0$ for some $l\in b$;
\item $x_i$ equals the point of $Q_i$ if $Q_i$ is a singleton.
\end{enumerate}
If every integer coordinate is recorded, let $P$ be the set of feasible
points that satisfy the recorded equalities and lie in $Q$. After the
substitutions described in the proof, let $Z$ be the tangent basis with
columns $e_i$ for free coordinates of blocks with a slack budget and
$e_i-e_{\iota}$ for the free coordinates other than an anchor $\iota$ of blocks
with an equality or recorded tight budget, let $c^P\in P$ be the center
constructed in the proof and $r^P$ the largest width of $Q$. If
\begin{equation}\label{eq:con:simplex-test}
 Z^{\mathsf T}\nabla^2F_\gamma(c^P)Z-(\kappa_3r^P+g_0)Z^{\mathsf T}Z\succ0,
\end{equation}
then $v^{\mathsf T}\nabla^2F_\gamma(x)v\ge g_0\norm v^2$ for $x\in P$ and $v$ in
the column space of $Z$, and the unique minimizer of $F_\gamma$ on $P$ is the
unique global minimizer. If the patch has tangent dimension zero, return
its rational point before constructing a center or testing a matrix.
Conversely, if $g_*\ge g_0$, every multiplier of an
active inequality at the minimizer $a$ exceeds $\tau$, $A=2+nL/g_0$, and
$h_j\le\min\{1/(4A),\tau/(8\kappa_2A),g_0/(8\kappa_3A)\}$, then every active
inequality at $a$ and every integer coordinate is recorded, and closure
returns a point or \eqref{eq:con:simplex-test} holds. Active inequalities
are zero-coordinate constraints and tight budgets of inequality blocks;
original equality budgets have no multiplier-margin requirement.
\end{lemma}

\begin{proof}
\emph{Soundness.} Every minimizer $x^*$ lies in $Q$, and
$\abs{\partial_iF_\gamma(x)-\partial_iF_\gamma(c)}\le\kappa_2r$ for $x\in Q$. In
(a), $\partial_iF_\gamma(x^*)>0$, and if $x_i^*>0$ then decreasing $x_i$ is
feasible in an inequality block and decreases the objective. In (b),
$\partial_lF_\gamma(x^*)<0$, and if the budget were slack, increasing $x_l$
would be feasible. In (c), $\partial_iF_\gamma(x^*)>\partial_lF_\gamma(x^*)$,
and if $x_i^*>0$, moving mass from $x_i$ to $x_l$ keeps every simplex
constraint and decreases the objective. These directions are feasible in the
original simplex; the hull is used only to bound derivatives. Case~(d) holds
because $Q$ contains every minimizer.

\emph{The patch.} In each block let $\mathsf F$ be the unrecorded coordinates
with bounds $l_i<u_i$ from $Q$ and $[0,1]$, and $\beta_b$ the budget left
after substituting fixed coordinates, with $\sum_{\mathsf F}x_i\le\beta_b$, or
$=\beta_b$ for an equality block or recorded tight budget. If
$\mathsf F=\emptyset$, drop this block before any division: containment of
a global minimizer certifies that its fixed values satisfy its budget.
If an equality budget has $\abs{\mathsf F}=1$, fix its remaining coordinate
to $\beta_b$ and drop the block. If
$\sum_{\mathsf F}l_i=\beta_b$, the block is forced to $l$; if the budget is an
equality and $\sum_{\mathsf F}u_i=\beta_b$, it is forced to $u$; substitute
such blocks. If no tangent direction remains, $P$ is a rational point and
closure returns it. Otherwise, every remaining block has a positive sum of
coordinate widths. For an equality budget put
$\lambda=(\beta_b-\sum l_i)/\sum(u_i-l_i)\in(0,1)$, and for an inequality
budget $\lambda=\min\{\frac12,(\beta_b-\sum l_i)/(2\sum(u_i-l_i))\}$, and let
$c^P_{\mathsf F}=l+\lambda(u-l)$. Every bound and every inequality budget is
strict at $c^P$, so the affine hull of $P$ is cut out exactly by the
recorded and equality budgets, and $Z$ spans its direction space. If no free
coordinate remains, $P$ is a single point, which is the unique minimizer.
For $x\in P$, $\norm{x-c^P}_\infty\le r^P$, so
$Z^{\mathsf T}\nabla^2F_\gamma(x)Z\succeq Z^{\mathsf T}\nabla^2F_\gamma(c^P)Z
-\kappa_3r^PZ^{\mathsf T}Z\succ g_0Z^{\mathsf T}Z$, which is the claimed bound
with $v=Zt$. The patch contains every minimizer, so its unique minimizer is
the unique global one.

\emph{Success.} The witnesses lie within $\frac{h_j}2\sqrt{nL/g_0}$ of $a$, so
every hull lies within $Ah_j$ of $a$: $r\le Ah_j$ and $r^P\le2Ah_j$. Integer
hulls are the correct singletons because $Ah_j\le1/4$. KKT conditions at $a$
read $\partial_iF_\gamma(a)+\lambda_b-\lambda_i=0$, with
$\lambda_i\ge0$ for zero-coordinate inequalities and
$\lambda_b\ge0$ for inequality budgets, while the multiplier
$\lambda_b$ of an original equality budget is unrestricted. They are
uniquely determined because the active normals, together with original
equality normals, are linearly independent. In a slack inequality block,
a zero coordinate has
$\partial_iF_\gamma(a)=\lambda_i>\tau$, and test (a) fires since
$\partial_iF_\gamma(c)-\kappa_2r\ge\lambda_i-2\kappa_2Ah_j>0$. In a block with a
tight inequality budget, a positive coordinate $l$ has
$\partial_lF_\gamma(a)=-\lambda_b<-\tau$, so test (b) fires. In every block
with a tight or equality budget, a zero coordinate $i$ has
$\partial_iF_\gamma(a)-\partial_lF_\gamma(a)=\lambda_i>\tau$, so test (c) fires
since its error is at most $4\kappa_2Ah_j\le\tau/2$. No lower bound on the
positive coordinate is needed. All original active inequalities have now
been recorded, and further singleton fixing is sound by containment. Hence
$P$ lies in the smallest face of $a$, the column space of $Z$ lies in the
direction space of that face, and two-sided growth gives
$Z^{\mathsf T}\nabla^2F_\gamma(a)Z\succeq2g_0Z^{\mathsf T}Z$. Since
$\norm{c^P-a}_\infty\le2Ah_j$, the matrix of \eqref{eq:con:simplex-test} is at
least $(g_0-4\kappa_3Ah_j)Z^{\mathsf T}Z\succeq\frac{g_0}2Z^{\mathsf T}Z$.
\end{proof}

\begin{lemma}[Multiplier margins on simplices]\label{lem:con:simplex-tail}
Let $D=\max\{1,d-1\}$ and
\[
 K_\Delta=\max\{1,\ 2n_cN_Z3^{n_c}(2D)^{n_c}\}.
\]
For $\tau>0$, the probability that $g_*>0$ and some active inequality at the
unique minimizer has multiplier at most $\tau$ is at most
$K_\Delta(\tau/\sigma+1/M)$. Only zero-coordinate inequalities and tight
budgets of inequality blocks enter this event; equality-budget
multipliers are excluded.
\end{lemma}

\begin{proof}
Fix integer labels $z$, a face of every block (support and tightness, at most
$3^{\abs b}$ choices), and one active inequality $\varkappa$ of the resulting
product face $\Phi$ (at most $2n_c$ choices). Let $t\in\R^k$ be the tangent
coordinates: the support coordinates of slack blocks and the support
coordinates other than the anchor $\iota$ (least support index) of tight
blocks. On $\Phi$ the linear noise equals a constant plus
$\theta^{\mathsf T}t$, where $\theta$ consists of the $\gamma_i$ of slack
supports and the differences $\gamma_i-\gamma_\iota$ of tight supports. The
stationarity system of $\Phi$, $\nabla_t[F_0(x(t))]+\theta=0$, has at most
$D^k$ nonsingular zeros by \cref{lem:sp:bezout} and depends on the noise only
through $\theta$. For $k=0$ there is one empty stationary tuple, counted
with weight $D^0=1$, without applying the positive-dimensional zero lemma.
At a point of $\Phi$ the active-inequality multipliers are
$\lambda_b=-\partial_\iota F_0-\gamma_\iota$ for a tight inequality budget,
$\lambda_i=\partial_iF_0-\partial_\iota F_0+\gamma_i-\gamma_\iota$ for a zero
coordinate of a tight block, and $\lambda_i=\partial_iF_0+\gamma_i$ for a zero
coordinate of a slack block. Thus $\lambda_\varkappa$ is affine with slope
$\pm1$ in its \emph{normal coefficient} $\nu_\varkappa$, namely $\gamma_\iota$
or $\gamma_i$, once $\theta$ and the other coefficients are fixed.

The map from $\gamma$ to $(\theta,\nu)$, where $\nu$ lists the original
coefficients not differenced, is injective. Its transformed coordinates are
not independent under the grid law, so we count tuples. A difference of two
grid values takes at most $2M-1$ values. For each of the at most
$(2M-1)^kM^{n-k-1}$ values of the transformed coordinates other than
$\nu_\varkappa$, the zeros are fixed, and for each zero the condition
$\abs{\lambda_\varkappa}\le\tau$ allows at most $\tau(M-1)/\sigma+1$ grid values
of $\nu_\varkappa$. Dividing by $M^n$, the probability that the event occurs
at some nonsingular zero is at most $2^kD^k(\tau/\sigma+1/M)$.

If $g_*>0$, let $a$ be the minimizer and $\Phi$ its smallest face. Its
tangent coordinates are a zero of the system, nonsingular because two-sided
growth along tangent directions gives a positive definite tangent Hessian.
The event is therefore contained in the union over $N_Z3^{n_c}2n_c$ choices
of the events above, with $k\le n_c$.
\end{proof}

\begin{proof}[Proof of \cref{thm:con:simplex}]
\emph{Schedule.} Let $B$ be the factor of \cref{thm:count:fallback} for the
product of simplices and integer intervals, and $C_{\rm tail}$ the constant of
\cref{lem:count:finite-tails} for its linear description and label
disjunctions. Put
\[
 \rho=\frac1{4B},\qquad g_0=\frac{\rho\sigma}{2W},\qquad
 \tau=\frac{\rho\sigma}{2K_\Delta},\qquad A=2+\frac{nL}{g_0},
\]
let $J$ be the least level with
$h_J\le\min\{1/(4A),\tau/(8\kappa_2A),g_0/(8\kappa_3A)\}$, and let $M$ be the
least power of two with $M\ge\max\{2,2^J,4nC_{\rm tail}/\rho,2K_\Delta/\rho\}$.
Every
coordinate width is $1$ for simplex coordinates and $w_i$ for integers, so
\cref{lem:count:finite-tails}(b) gives $\Prob\{g_*<g_0\}\le\rho$, and
\cref{lem:con:simplex-tail} gives probability at most $\rho$ for small
multipliers. By \cref{lem:con:simplex-close} closure succeeds by level $J$
outside these events, and is sound on every draw; otherwise the fallback
runs on the same draw.

\emph{Work.} The cells have at most $2^{\abs b}$ corners, children and
incident cells per block; \cref{lem:con:simplex-count} and the argument of
\cref{thm:sp:main} give the expected work, and the fallback contributes
$\poly_d(I)$ by \cref{lem:count:rare-fallback}(a).

\emph{Evaluation.} If the tangent dimension is zero, return the rational
point and its exactly evaluated objective; no convex-body oracle is used.
Otherwise apply \cref{lem:sp:gls} to the patch $P$ in the free
coordinates, with $c=c^P$, the basis $Z$, $\zeta=2k$ (since
$\norm Z_2\le\norm Z_F\le\sqrt{2k}$), $R=\max\{1,k\}$, and
$r=\min\{\frac12,\min_s\sigma_s/(k\norm{\alpha_s}_1)\}$, where $\alpha_s^{\mathsf
T}x\le\alpha_{0,s}$ runs over the bounds and inequality budgets of $P$ and
$\sigma_s>0$ is its slack at $c^P$: every coordinate of $Zt$ is a coordinate of
$t$ or minus a sum of them, so $\abs{\alpha_s^{\mathsf T}Zt}\le
k\norm{\alpha_s}_1\norm t_2$. The repair map is the Euclidean projection onto
$P$, computed blockwise: on a block with bounds $l\le x\le u$ and budget
$\underline\beta\le\sum x\le\beta$ (with $\underline\beta=\beta$ for equality
budgets and $\underline\beta=-\infty$ otherwise), the projection of $w$ is
$x_i=\min\{u_i,\max\{l_i,w_i-\theta\}\}$, with $\theta=0$ if this satisfies
the budget and otherwise the solution of $\sum_ix_i(\theta)$ equal to the
violated endpoint. The sum is continuous, nonincreasing and piecewise linear
in $\theta$ with breakpoints $w_i-u_i$ and $w_i-l_i$; sorting them and solving
one linear equation gives a rational $\theta$ in polynomial time. Projection
onto a convex set is nonexpansive, so $\nu=\norm\cdot_2$. The curvature
hypothesis holds with $g=g_0$ by \cref{lem:con:simplex-close}. Use objective
tolerance $\min\{2^{-q},(g_0/2)2^{-2q}\}$, so the lemma controls physical
Euclidean distance and the feasible objective gap. On fallback draws, let
$G'\ge\max\{1,\sup\norm{\nabla F_\gamma}_2\}$ bound the gradient on the
enclosing box. Refine each continuous coordinate to error at most
$2^{-q-1}/(nG')$, project each block onto its original simplex in the same
way, and keep integer labels exact. Nonexpansiveness gives physical distance
at most $2^{-q-1}/G'$ and objective gap at most $2^{-q-1}$. A value enclosure
of width $2^{-q-1}$ then completes the certified interval. These refinements
cost $B\poly_d(I+q)$ and have expectation $\poly_d(I+q)$.
\end{proof}

\subsection{Order constraints}\label{app:con:order}

Write $x=(x_C,x_Z)$ for continuous and binary coordinates. The slice for a
binary assignment $\mathsf b$ is
$P_{\mathsf b}=\{x_C:(x_C,\mathsf b)\in X\}$, an order polytope on $I_C$
intersected with the bounds $x_i=0$ for edges $(i,l)$ with $\mathsf b_l=0$ and
$x_i=1$ for edges $(l,i)$ with $\mathsf b_l=1$.

\begin{lemma}[Zero-one vertices]\label{lem:con:order-vertices}
Every vertex of $P_{\mathsf b}$, and of every face of it, is a zero-one vector.
\end{lemma}

\begin{proof}
Let $x$ be a point of $P_{\mathsf b}$ with a value $t\in(0,1)$, and let $G$ be
the set of coordinates with value $t$. Moving all coordinates of $G$ by
$\pm\epsilon$ keeps every edge inside $G$ tight, keeps every edge between
$G$ and other coordinates strict for small $\epsilon$, and keeps the bounds;
so $x$ is the midpoint of two feasible points and not a vertex. A face is a
polytope whose vertices are vertices of $P_{\mathsf b}$.
\end{proof}

\begin{proof}[Rounding for \cref{thm:con:order}]
For fixed $U\in[0,1]$ the scalar map
$t\mapsto h(\lfloor t/h\rfloor+\ind\{U<\operatorname{frac}(t/h)\})$ is
nondecreasing, fixes the grid values, including $0$ and $1$, and maps each
grid interval onto its endpoints. Applied to all continuous coordinates of
$x\in\Omega_j$, it therefore preserves every order constraint and every
binary value, keeps every coordinate strictly inside a grid interval within
that interval and every grid coordinate fixed, and so keeps every bag tuple
a corner of every candidate cell containing it. For $U$ uniform, every
coordinate is mean preserving with variance at most $h^2/4$. The segment from
$x$ to $Y$ lies in $[0,1]^n$, so the full Hessian bound gives
$\E F_\gamma(Y)\le F_\gamma(x)+\frac{\bar L}2\E\norm{Y-x}^2\le F_\gamma(x)+n_c\bar
Lh^2/8$. Every edge lies in a bag and rows violating an assigned edge are
removed, so allowed grid points are feasible. The proof of
\cref{prop:sp:prune} then applies with $E_j=n_c\bar Lh_j^2/8$.
\end{proof}

Fix a bag $\beta$, write $\beta_C=\beta\cap I_C$ with $c_\beta=\abs{\beta_C}$ and
$\beta_Z=\beta\cap I_Z$, and let $O=[n]\setminus\beta$. For a binary row
$\mathsf b\in\{0,1\}^{\beta_Z}$ and $v\in[0,1]^{\beta_C}$, put
\[
 V_{\mathsf b}(v)=\gamma_{\beta_Z}^{\mathsf T}\mathsf b+\min\bigl\{F_0(v,\mathsf
 b,z)+\gamma_O^{\mathsf T}z:\ (v,\mathsf b,z)\in X\bigr\},
\]
with $\min\emptyset=+\infty$; it depends on $\gamma_O$ and $\gamma_{\beta_Z}$
only. A \emph{standard order simplex} is
$S_\pi=\{v:0\le v_{\pi(1)}\le\dots\le v_{\pi(c_\beta)}\le1\}$ for a permutation
$\pi$ of $\beta_C$. Its vertices form a chain of zero-one vectors, and a face
with vertices $\mathsf a_0<\dots<\mathsf a_k$ in this chain has \emph{edge
directions} $\mathsf d_l=\mathsf a_l-\mathsf a_{l-1}$, $l=1,\dots,k$, with
disjoint nonempty supports.

\begin{lemma}[Endpoint-preserving transport]\label{lem:con:transport}
Let $v$ lie in the relative interior of a face $\Phi$ of $S_\pi$, with
$V_{\mathsf b}(v)<\infty$. There is an affine map $w\mapsto x(w)$ on the
affine hull of $\Phi$ such that, for every $w\in\Phi$, $x(w)\in X$ has
continuous bag part $w$ and binary bag part $\mathsf b$; $x(v)$ attains
$V_{\mathsf b}(v)$; and along every edge direction $\mathsf d_l$ of $\Phi$ each
coordinate of $x$ changes at a rate in $[0,1]$, binary coordinates at rate
$0$. Consequently, whenever $v\pm\theta\mathsf d_l\in\Phi$,
\[
 V_{\mathsf b}(v+\theta\mathsf d_l)+V_{\mathsf b}(v-\theta\mathsf d_l)
 -2V_{\mathsf b}(v)\le n_c\bar L\theta^2 .
\]
\end{lemma}

\begin{proof}
On the relative interior of $\Phi$, coordinates of $\beta_C$ are partitioned
into groups of equal value: a group at $0$, a group at $1$ (either may be
empty), and $k$ interior groups with values $0<r_1<\dots<r_k<1$; this
partition is the same at every relative-interior point, and on the closed
face the groups stay equal and ordered, $0\le s_1(w)\le\dots\le s_k(w)\le1$,
where $s_l(w)$ is the common value of group $l$ at $w$, a linear function of
$w$. Moving along $\mathsf d_l$ changes exactly one group value at unit rate:
writing $v=\sum_l\lambda_l\mathsf a_l$ with positive weights, the value of
the $l$th interior group is $\sum_{i\ge l}\lambda_i$, and $\mathsf d_l$ moves
weight from $\lambda_{l-1}$ to $\lambda_l$.

Let $x^\star=(v,\mathsf b,z^\star)$ attain $V_{\mathsf b}(v)$; it exists by
compactness. Put $r_0=0$, $r_{k+1}=1$, $s_0=0$, $s_{k+1}=1$, and let $T_w$ be
the piecewise-affine map of $[0,1]$ with $T_w(r_l)=s_l(w)$, affine between
consecutive knots. Define $x(w)$ by applying $T_w$ to every coordinate of
$x^\star$. For $w\in\Phi$, $T_w$ is nondecreasing with $T_w(0)=0$ and
$T_w(1)=1$, so $x(w)$ satisfies every order constraint, keeps every binary
value, lies in $[0,1]^n$, and maps every bag coordinate of group $l$ to
$s_l(w)=w_i$; thus $x(w)\in X$ with bag parts $(w,\mathsf b)$, and $x(v)=x^\star$.
A coordinate of $x^\star$ in $[r_l,r_{l+1}]$ equals
$(1-\vartheta)r_l+\vartheta r_{l+1}$ for a fixed $\vartheta\in[0,1]$ and is
mapped to $(1-\vartheta)s_l(w)+\vartheta s_{l+1}(w)$, which is affine in $w$;
along $\mathsf d_l$ its rate is $1-\vartheta$, $\vartheta$, or $0$. Binary
coordinates sit at the fixed knots $0$ and $1$.

The function $\varphi(w)=\gamma_{\beta_Z}^{\mathsf T}\mathsf b+F_0(x(w))+
\gamma_O^{\mathsf T}x(w)_O$ satisfies $\varphi(v)=V_{\mathsf b}(v)$ and
$\varphi\ge V_{\mathsf b}$ on $\Phi$. The linear noise term is affine in $w$,
and the derivative $\mathsf x'$ of $x$ along $\mathsf d_l$ has at most $n_c$
nonzero entries, all in $[0,1]$, so
$\frac{d^2}{d\theta^2}\varphi(v+\theta\mathsf d_l)=\mathsf x'^{\mathsf
T}\nabla^2F_0\mathsf x'\le\bar L\norm{\mathsf x'}^2\le n_c\bar L$. Comparing
$V_{\mathsf b}$ with its upper support $\varphi$ at $v\pm\theta\mathsf d_l$
gives the inequality.
\end{proof}

The transported points need not exhaust the fiber over $w$: at endpoint faces
new binary labels can become feasible, and the projection of the mixed domain
onto the continuous bag coordinates need not be convex. Neither matters,
because only the upper support at $v\pm\theta\mathsf d_l$ is used.

\begin{proposition}[Order count]\label{prop:con:order-count}
Let $h=h_j=2^{-j}$ with $Mh\ge1$, and let $N_j(\beta)$ be the set of feasible
bag grid tuples $(v,\mathsf b)$ with
$V_{\mathsf b}(v)+\gamma_{\beta_C}^{\mathsf T}v\le f^*+2E_j$. Then
\[
 \E\abs{N_j(\beta)}\le2^{\abs{\beta_Z}}c_\beta!\,(c_\beta+1)
 \Bigl[2+\frac{3n_c\bar L}{4\sigma}\Bigr]^{c_\beta},
\]
and $\abs{\mathcal W'_j(\beta)}\le2^{\abs\beta}\abs{N_j(\beta)}$ on every draw.
\end{proposition}

\begin{proof}
Fix $\mathsf b$ and a feasible $v$, choose $\pi$ with $v\in S_\pi$, and let
$\Phi$ be the face of $S_\pi$ whose relative interior contains $v$, of
dimension $k$. Its barycentric weights are positive multiples of $h$, so
$v\pm h\mathsf d_l\in\Phi$ for every edge direction. Condition on all
coefficients except those of $\beta_C$. If $(v,\mathsf b)\in N_j(\beta)$, the
argument of \cref{lem:sp:compare}, with the transported points of
\cref{lem:con:transport} as feasible competitors and $\eta=2E_j=n_c\bar Lh^2/4$,
confines $\gamma^{\mathsf T}\mathsf d_l$ to an interval of length at most
$(n_c\bar Lh^2+2\eta)/h=\frac32n_c\bar Lh$ that is determined by the
conditioned values and the tuple. Choose one coordinate in the support of
each $\mathsf d_l$ and condition also on the other continuous bag
coefficients; the supports are disjoint, so the $k$ remaining coefficients
are independent and each lies in an interval of that length. The conditional
probability is at most $[h(\frac{3n_c\bar L}{4\sigma}+\frac1{Mh})]^k$. A
$k$-face has $\binom{1/h-1}{k}\le h^{-k}/k!$ relative-interior grid points,
$S_\pi$ has $\binom{c_\beta+1}{k+1}$ faces of dimension $k$, there are at most
$c_\beta!$ permutations, and at most $2^{\abs{\beta_Z}}$ binary rows; tuples
lying in several simplices are counted more than once, which only enlarges
the bound. With $\mathsf A=\frac{3n_c\bar L}{4\sigma}+1\ge\frac{3n_c\bar
L}{4\sigma}+\frac1{Mh}$,
\[
 \E\abs{N_j(\beta)}\le2^{\abs{\beta_Z}}c_\beta!\sum_{k=0}^{c_\beta}
 \binom{c_\beta+1}{k+1}\frac{\mathsf A^k}{k!}
 \le2^{\abs{\beta_Z}}c_\beta!\,(c_\beta+1)(1+\mathsf A)^{c_\beta},
\]
using $\binom{c+1}{k+1}/k!=(c+1)\binom ck/(k+1)!\le(c+1)\binom ck$. The second
claim follows from \cref{prop:sp:prune}(d), since a retained cell has an
allowed witness $y$ with $V_{\mathsf b}(v)+\gamma_{\beta_C}^{\mathsf T}v\le
F_\gamma(y)$.
\end{proof}

\begin{lemma}[Linear-programming exposure]\label{lem:con:lpgap}
Suppose every global minimizer has binary part $\mathsf b$, let $Q_C$ be a box
containing their continuous parts, $c$ its midpoint, $r$ its largest
half-width, $g=\nabla_CF_\gamma(c,\mathsf b)$ and $\delta=\kappa_2r$, where
$\kappa_2$ bounds the continuous Hessian row sums on $[0,1]^n$. For a proposed
equality $\mathsf e$ (an edge $x_i=x_l$ with opposite face
$P_{\mathsf e}^-=P_{\mathsf b}\cap\{x_i=0,x_l=1\}$, or a bound $x_i=0$ or $x_i=1$
with opposite face $P_{\mathsf b}\cap\{x_i=1\}$ or $P_{\mathsf b}\cap\{x_i=0\}$)
with nonempty opposite face, put
$\Delta_{\mathsf e}(v)=\min_{P_{\mathsf e}^-}v^{\mathsf T}x-\min_{P_{\mathsf b}}v^{\mathsf T}x$.
If $\Delta_{\mathsf e}(g)>n_c\delta$, the equality holds at every global
minimizer. For all $v,w$, $\abs{\Delta_{\mathsf e}(v)-\Delta_{\mathsf e}(w)}\le
n_c\norm{v-w}_\infty$. If the opposite face is empty, the equality holds on all
of $P_{\mathsf b}$.
\end{lemma}

\begin{proof}
Let $x^*=(a,\mathsf b)$ be a global minimizer and $\mathsf h=\nabla_CF_\gamma(a,\mathsf
b)$. Since $a$ minimizes $F_\gamma(\cdot,\mathsf b)$ on the convex set
$P_{\mathsf b}$, $\mathsf h^{\mathsf T}(x-a)\ge0$ there, so $a$ lies in the face
$\Psi=\argmin_{P_{\mathsf b}}\mathsf h^{\mathsf T}x$, whose vertices are zero-one
by \cref{lem:con:order-vertices}. If an edge equality fails at $a$, then
$a_i<a_l$, and not all vertices of $\Psi$ have equal $i$th and $l$th
coordinates; a vertex $\mathsf v$ with $\mathsf v_i\ne\mathsf v_l$ has
$\mathsf v_i=0$ and $\mathsf v_l=1$, so $\mathsf v\in P_{\mathsf e}^-$. For a
bound, a vertex of $\Psi$ with the opposite value exists likewise. Let
$\mathsf w$ be a zero-one minimizer of $g^{\mathsf T}x$ on $P_{\mathsf b}$. Since
$\mathsf h^{\mathsf T}\mathsf v\le\mathsf h^{\mathsf T}\mathsf w$,
$\norm{\mathsf v-\mathsf w}_1\le n_c$, and $\norm{g-\mathsf h}_\infty\le\kappa_2
\norm{c-a}_\infty\le\delta$,
\[
 \Delta_{\mathsf e}(g)\le g^{\mathsf T}(\mathsf v-\mathsf w)
 =\mathsf h^{\mathsf T}(\mathsf v-\mathsf w)+(g-\mathsf h)^{\mathsf T}(\mathsf v-\mathsf w)
 \le n_c\delta ,
\]
a contradiction. For the Lipschitz bound, let $\mathsf x_1$ be a zero-one
minimizer of $w^{\mathsf T}x$ on $P_{\mathsf e}^-$ and $\mathsf x_2$ one of
$v^{\mathsf T}x$ on $P_{\mathsf b}$; then
$\Delta_{\mathsf e}(v)-\Delta_{\mathsf e}(w)\le(v-w)^{\mathsf T}(\mathsf x_1-\mathsf
x_2)\le n_c\norm{v-w}_\infty$, and symmetrically. If $P_{\mathsf e}^-$ is empty,
no vertex of $P_{\mathsf b}$ violates the equality, so it holds on the convex
hull of the vertices.
\end{proof}

\paragraph{The closure test.} After retention at level $j\ge1$: form the hulls;
fail unless every binary hull is a singleton, which fixes $\mathsf b$; apply
\cref{lem:con:lpgap} to every continuous edge and every continuous bound,
using rational linear programs; merge coordinates joined by certified edge
equalities into blocks, fix blocks with certified bounds or bounds forced by
$\mathsf b$, intersect the hull intervals of the members, contract directed
cycles, and propagate bounds, replacing each lower bound by the largest lower
bound of its predecessors and each upper bound by the smallest upper bound of
its successors; fix blocks whose interval is a single point and repeat until
nothing changes; remove tautological rows. This yields $x_C=Dy+\bar x_C$
with disjoint zero-one columns, a box $[\mathsf L,\mathsf U]$ in $y$ with
$\mathsf L<\mathsf U$, and an acyclic set of remaining order inequalities
$y_u\le y_{u'}$. If no block remains, return the point. Otherwise, with $c_y$
the midpoint and $r_y$ the largest half-width of the box, test
\begin{equation}\label{eq:con:order-test}
 D^{\mathsf T}\nabla^2_{CC}F_\gamma(Dc_y+\bar x_C,\mathsf b)D
 -(\kappa_3r_y+g_0)D^{\mathsf T}D\succ0,
\end{equation}
where $\kappa_3$ bounds the continuous third-derivative row sums on
$[0,1]^n$.

\begin{lemma}[Order closure]\label{lem:con:order-close}
The patch $P_y=\{y\in[\mathsf L,\mathsf U]:y_u\le y_{u'}\text{ for the remaining
inequalities}\}$ contains the block values of every global minimizer and is
full-dimensional. If \eqref{eq:con:order-test} holds, then
$\nabla^2_yF_\gamma(Dy+\bar x_C,\mathsf b)\succeq g_0D^{\mathsf T}D\succeq g_0I$ on
$P_y$, and the unique minimizer on $P_y$ is the unique global minimizer.
Conversely, if $g_*\ge g_0$, every active proposal at the minimizer with a
nonempty opposite face has exposure gap larger than $\tau$, $A=2+n_c\bar
L/g_0$, and $h_j\le\min\{1/(4A),\tau/(4n_c\kappa_2A),g_0/(4\kappa_3A)\}$, then the
test of level $j$ succeeds.
\end{lemma}

\begin{proof}
\emph{Containment.} Binary singleton hulls fix $\mathsf b$ at every minimizer,
the certified equalities hold by \cref{lem:con:lpgap}, and merging, cycle
contraction, bound propagation and singleton fixing are equivalent
reformulations of $P_{\mathsf b}\cap Q_C$ with those equalities.

\emph{Full dimension.} After propagation the bounds are monotone along the
remaining inequalities: $\mathsf L_u\le\mathsf L_{u'}$ and $\mathsf U_u\le\mathsf
U_{u'}$ for $u\to u'$. For an inequality $u\to u'$, the point that equals
$\mathsf U$ on the blocks reachable from $u'$ and $\mathsf L$ elsewhere is
feasible, and $u$ is not reachable from $u'$ because the remaining graph is
acyclic: an inequality from an unreachable block $w$ to a reachable $w'$
has $\mathsf L_w\le\mathsf L_{w'}\le\mathsf U_{w'}$, and none goes from a
reachable to an unreachable block. At this point $y_u=\mathsf L_u$ and
$y_{u'}=\mathsf U_{u'}$. If $\mathsf L_u=\mathsf U_{u'}$, then
$\mathsf L_u\le\mathsf L_{u'}\le\mathsf U_{u'}$ and
$\mathsf L_u\le\mathsf U_u\le\mathsf U_{u'}$ force singleton intervals, which
were removed. So no inequality and, by the points $\mathsf L$ and $\mathsf U$,
no bound is tight on all of $P_y$, and $P_y$ has nonempty interior.

\emph{Certificate.} For $y,y'$ in the box,
$\norm{D(y-y')}_\infty=\norm{y-y'}_\infty$ and
$D^{\mathsf T}(\nabla^2F_\gamma(x)-\nabla^2F_\gamma(x'))D\succeq-\kappa_3\norm{x-x'}_\infty
D^{\mathsf T}D$, so \eqref{eq:con:order-test} gives the claimed bound, and
$D^{\mathsf T}D\succeq I$ because every block has at least one coordinate.

\emph{Success.} Witnesses lie within $\frac{h_j}2\sqrt{n_c\bar L/g_0}$ of the
minimizer $x^*=(a,\mathsf b^*)$, so hulls lie within $Ah_j$ of it, and binary
hulls are the singletons $\mathsf b^*$ because $Ah_j\le1/4$ and binary cells
are singletons at levels $j\ge1$. By \cref{lem:con:lpgap},
$\Delta_{\mathsf e}(g)\ge\Delta_{\mathsf e}(\mathsf h)-n_c\kappa_2r>\tau-\tau/4
>n_c\delta$ for every active proposal with a finite gap, so all of them are
certified. Proposals with an empty opposite face hold on all of
$P_{\mathsf b}$. Hence every point of $P_y$ satisfies every constraint that is
tight at $a$, the patch lies in the smallest face of $a$ in $P_{\mathsf b}$,
and the columns of $D$ times directions of $P_y$ are directions of that face.
Two-sided growth along them gives
$D^{\mathsf T}\nabla^2F_\gamma(x^*)D\succeq2g_0D^{\mathsf T}D$. With
$\norm{Dc_y+\bar x_C-a}_\infty\le r_y\le Ah_j$, the matrix in
\eqref{eq:con:order-test} is at least
$(g_0-2\kappa_3Ah_j)D^{\mathsf T}D\succeq\frac{g_0}2D^{\mathsf T}D$.
\end{proof}

This weighted test measures curvature in the metric of the original
coordinates. An unweighted version that subtracts $(n_c\kappa_3r_y+g_0)I$ is
also sound, but the two must not be mixed by dropping the factor $D^{\mathsf
T}D$ from one and the factor $n_c$ from the other.

\begin{lemma}[Exposure gaps]\label{lem:con:order-tail}
Let $D=\max\{1,d-1\}$ and
$K_\preceq=\max\{1,2^{n_z}(m+2n_c)2^{m+2n_c}D^{n_c}\}$. For $\tau>0$, the
probability that $g_*>0$ and some proposal that is active at the minimizer and
has a nonempty opposite face has exposure gap
$\Delta_{\mathsf e}(\nabla_CF_\gamma(x^*))\le\tau$ is at most
$K_\preceq(\tau/\sigma+2/M)$.
\end{lemma}

\begin{proof}
Fix a binary assignment $\mathsf b$, a face $\Phi$ of $P_{\mathsf b}$ (determined
by a subset of its at most $m+2n_c$ inequalities), and a proposal $\mathsf e$
that holds on $\Phi$. Write the affine hull of $\Phi$ as $x=x_\Phi+B\zeta$ with
$B$ of full column rank. Its stationarity system
$B^{\mathsf T}(\nabla_CF_0(x_\Phi+B\zeta,\mathsf b)+\gamma_C)=0$ has at most
$D^{n_c}$ nonsingular zeros by \cref{lem:sp:bezout}. If $\Phi$ has dimension
zero there is one empty stationary tuple, counted with weight one.

For an edge $\mathsf e=(i,l)$, condition on all coefficients except
$\gamma_i,\gamma_l$, and write $\gamma_i=t$, $\gamma_l=s-t$. Every direction of
$\Phi$ has equal $i$th and $l$th entries, so $B^{\mathsf T}(e_i-e_l)=0$ and the
system depends on $s$ but not on $t$. Fix $s$ and a zero, with point
$x_\rho$ in the relative interior of $\Phi$. Whenever $(x_\rho,\mathsf b)$ is a
global minimizer with smallest face $\Phi$, the slice gradient is
$\mathsf h(t)=\mathsf h(0)+t(e_i-e_l)$, every vertex $\mathsf w$ of $\Phi$
minimizes $\mathsf h(t)^{\mathsf T}x$ on $P_{\mathsf b}$ (the minimizing face contains
$x_\rho$ and hence $\Phi$), and every point of $P_{\mathsf e}^-$ has $x_i=0$,
$x_l=1$, while $\mathsf w_i=\mathsf w_l$. So
$\Delta_{\mathsf e}(\mathsf h(t))=\beta_\rho-t$ with $\beta_\rho$ independent of
$t$, and the event $\Delta_{\mathsf e}\le\tau$ confines $t$ to an interval of
length $\tau$, since gaps are nonnegative. For each of the at most $2M-1$
values of $s$, the admissible values of $t$ are spaced by $2\sigma/(M-1)$, so
at most $\tau(M-1)/(2\sigma)+1$ of them lie in the interval of each zero of
that fiber. Summing over the fibers and dividing by $M^2$, the probability for
one face and edge is at most
$D^{n_c}(2M-1)[\tau(M-1)/(2\sigma)+1]/M^2\le D^{n_c}(\tau/\sigma+2/M)$. A tuple-count on the
fibers of $s$ is used here because the conditional law of $t$ given $s$ is not
uniformly spread. For a bound $x_i=0$, vary $\gamma_i$ alone: $B^{\mathsf
T}e_i=0$, and $\Delta_{\mathsf e}=\beta_\rho+\gamma_i$; the upper bound is
symmetric; the probability is at most $\tau/(2\sigma)+1/M$.

If $g_*>0$, the minimizer $(a,\mathsf b)$ has a smallest face $\Phi$ in
$P_{\mathsf b}$, its zero is nonsingular because two-sided growth gives
$B^{\mathsf T}\nabla^2_{CC}F_\gamma B\succeq2g_*B^{\mathsf T}B\succ0$, and the
active proposals are those holding on $\Phi$. The union over at most
$2^{n_z}\cdot2^{m+2n_c}\cdot(m+2n_c)$ choices gives the bound. Edges between a continuous and a binary coordinate are bounds of the
slice and are covered; binary coefficients are fixed throughout.
\end{proof}

\begin{proof}[Proof of \cref{thm:con:order}]
\emph{Schedule.} Let $B$ be the factor of \cref{thm:count:fallback} for the
order polytope with binary constraints, $C_{\rm tail}$ the constant of
\cref{lem:count:finite-tails} for its description, $\rho=1/(4B)$,
$g_0=\rho\sigma/(2n)$, $\tau=\rho\sigma/(2K_\preceq)$, $A=2+n_c\bar L/g_0$, $J\ge1$
the least level satisfying the condition of \cref{lem:con:order-close}, and
$M$ the least power of two with
$M\ge\max\{2,2^J,4nC_{\rm tail}/\rho,4K_\preceq/\rho\}$. All widths are at most
one, so \cref{lem:count:finite-tails}(b) gives $\Prob\{g_*<g_0\}\le\rho$, and
\cref{lem:con:order-tail} gives probability at most
$K_\preceq\tau/\sigma+2K_\preceq/M\le\rho$ for small gaps. Outside these events
closure succeeds by level $J$; it is sound on every draw; otherwise the
fallback runs on the same draw.

\emph{Work.} Binary cells are $\{0,1\}$ at level $0$ and singletons
afterwards, so cells have at most $2^{\abs\beta}$ children and corners and a
grid tuple lies in at most $2^{\abs\beta}$ cells. \Cref{prop:con:order-count}
with $Mh_j\ge1$ for $j\le J$, and the argument of \cref{thm:sp:main}, give
\eqref{eq:con:orderbound}; the linear programs and matrix tests of a level
cost $\poly_d(I)$, and the fallback contributes $\poly_d(I)$ in expectation.

\emph{Evaluation.} If no block remains, evaluate the returned rational
point and its sampled value exactly. Otherwise apply \cref{lem:sp:gls} in the block coordinates $y$, with
$Z=I$ and $\zeta=1$. A rational center and inner radius come from the linear
program $\max\{\vartheta:\alpha_s^{\mathsf T}y+\vartheta\norm{\alpha_s}_1\le
\alpha_{0,s}\text{ for all rows},\ \vartheta\le\frac12\}$ over the rows of $P_y$:
its optimal vertex $(y_c,\vartheta^*)$ is rational of polynomial length, and
$\vartheta^*>0$ by full dimensionality; the ball of radius $\vartheta^*$ about
$y_c$ lies in $P_y$ since $\abs{\alpha_s^{\mathsf T}u}\le\norm{\alpha_s}_1\norm
u_2$. Take $R=\max\{1,\dim y\}$. The repair map clips to $[\mathsf L,\mathsf U]$
and then replaces every coordinate by the largest clipped value over its
predecessors, itself included. The result is feasible, because propagated
upper bounds are monotone along the inequalities. If $\bar y\in P_y$ and
$\norm{z-\bar y}_\infty=\epsilon$, clipping keeps every coordinate within
$\epsilon$ of $\bar y$, the repaired value is at least the clipped value and
hence at least $\bar y_u-\epsilon$, and it equals a clipped predecessor value
at most $\bar y_w+\epsilon\le\bar y_u+\epsilon$; so the repair is
nonexpansive in $\norm\cdot_\infty$, and we take $\nu=\norm\cdot_\infty$ with a
monomial bound on $\norm{D^{\mathsf T}\nabla_CF_\gamma}_1\le\norm{\nabla_CF_\gamma}_1$.
The curvature hypothesis holds with $g=g_0$ by \cref{lem:con:order-close}.
Since $\norm D_2\le\sqrt{n_c}$, the objective tolerance
$\min\{2^{-q},(g_0/(2n_c))2^{-2q}\}$ gives distance $2^{-q}$ in the original
coordinates. On fallback draws, identify the binary labels by refining to
error below $\frac12$, propagate their bounds through the order graph, clip the
continuous approximations to the propagated intervals, and take predecessor
maxima; this keeps the labels and every order constraint and does not
increase the $\infty$-norm error. Choose each continuous coordinate's error
at most $2^{-q-1}/(nG')$, where
$G'\ge\max\{1,\sup_{[0,1]^n}\norm{\nabla_CF_\gamma}_1\}$. Then the physical
Euclidean distance is at most $2^{-q-1}$ and the feasible objective gap at
most $2^{-q-1}$. Refining the optimal value to width $2^{-q-1}$ supplies the
certified value interval, in $B\poly_d(I+q)$ bit work. Without propagation,
a slightly positive approximation of a continuous predecessor could move a
binary zero.
\end{proof}
